% =====================================================================
% Self-summary (for overseer agent):
% This related-work section positions TensorGuard within six PL/ML
% traditions: refinement and dependent types (Liquid Haskell, F*, Dex,
% Hasktorch); rely/guarantee and contracts (Jones, Meyer, Findler--
% Felleisen); abstract interpretation for ML (Cousot--Cousot, Ariadne,
% DeepPoly, AI$^2$, jaxtyping); gradual/optional typing (Siek--Taha,
% MyPy, TypeScript); PyTorch- and autograd-specific tooling
% (FakeTensorMode, torch.export, TorchDynamo, Pytea, Hattori et al.);
% and lightweight static analysis (Pythia, ruff, beartype). Pytea is
% treated as the closest comparator and contrasted in detail. The
% contribution table (H_contribution_table.tex) is calibrated --- it
% marks the dimensions on which TensorGuard genuinely advances
% (refinement-typed grad-flag tracking, a Lean-checked core, abstention
% as a first-class verdict, no-execution backward verification on modern
% Transformer code) while marking ($\partial$/$\times$) the
% places where TG overlaps with or trails prior work (e.g.\ jaxtyping's
% lightweight ergonomics, Pytea's mature symbolic-shape engine,
% torch.export's industrial integration). Tone is calibrated and
% non-dismissive throughout. Citations marked TODO are uncertain venue/
% year combinations the overseer should verify.
% =====================================================================

\section{Related Work}
\label{sec:related}

TensorGuard sits at the intersection of several programming-languages
traditions that have, individually, addressed pieces of the
shape/dtype/grad-correctness problem. We organize the discussion by
tradition and conclude with a calibrated comparison
(Table~\ref{tab:contrib}).

\subsection{Refinement and dependent types}

The closest theoretical antecedents are refinement-type systems.
Liquid Haskell~\citep{vazou2014liquidhaskell,vazou2017refinement}
demonstrated that SMT-decidable refinements can scale to a real
programming language with minimal annotation burden; Liquid
Java~\citep{rondon2012liquidjava} and Sage~\citep{flanagan2006sage}
extended the approach to imperative and dynamically-typed settings.
Full dependent types in F$^\star$~\citep{swamy2016dependent},
Dafny~\citep{leino2010dafny}, Coq, and Lean~\citep{moura2021lean4}
allow shape constraints to be stated and proved as ordinary
propositions, but require human proof effort that is impractical for
day-to-day model code. Dependent
Haskell~\citep{eisenberg2016dependent} and the
\texttt{singletons}~\citep{eisenberg2012singletons} library brought
type-level naturals and singleton-indexed tensors to a mainstream
language, and Hasktorch~\citep{hasktorch} and the
\texttt{singletons-tensor} ecosystem applied them to deep learning;
these systems achieve elegant static guarantees but are confined to
Haskell-style codebases.

A parallel line of work targets array languages directly.
Dex~\citep{paszke2019dex,maclaurin2019dex} introduces typed indices
and effect tracking for differentiable array programming; \emph{Tensors
Fitting Perfectly}~\citep{rink2018tensors} formalizes a calculus where
shape mismatches are statically impossible. TensorGuard borrows the
refinement-type \emph{discipline} of these systems but applies it in
situ to unmodified PyTorch source rather than requiring a new host
language.

\subsection{Assume/guarantee reasoning and contracts}

The verdict structure of TensorGuard --- \textsc{Verified},
\textsc{Refuted}, \textsc{Abstain} --- inherits from
rely/guarantee~\citep{jones1983ar} and design-by-
contract~\citep{meyer1992contracts}. Higher-order
contracts~\citep{findler2002higherorder} and the
Misra--Chandy~\citep{misra1981proofs} treatment of distributed
specifications informed our handling of callable modules, where the
forward and backward contracts of a submodule must compose with those
of its caller. Unlike classical contract systems, TensorGuard's
contracts are \emph{checked statically without execution}, and
abstention is a first-class verdict rather than a fallback to
runtime checking.

\subsection{Abstract interpretation for ML}

Abstract interpretation~\citep{cousot1977ai} underlies most modern
robustness and shape verifiers for neural networks. Ariadne for
TensorFlow~\citep{dolby2018ariadne} pioneered abstract shape
inference for dataflow graphs.
DeepPoly~\citep{singh2019deeppoly} and
AI$^2$~\citep{gehr2018ai2} apply numerical abstract domains to
robustness certification --- a complementary problem to ours, since
they assume shapes are correct and reason about \emph{values}. On
the shape side, JAX's \texttt{jax.eval\_shape} and
\texttt{jax.make\_jaxpr} provide concrete shape propagation under
tracing, and \texttt{jaxtyping}~\citep{jaxtyping}/
\texttt{torchtyping}~\citep{torchtyping} expose shape annotations as
runtime-checked Python types. TensorGuard differs by performing
\emph{symbolic} shape reasoning ahead of any tracer execution and by
modeling \texttt{requires\_grad} as part of the refinement.

\subsection{Gradual and optional typing}

Siek and Taha's gradual typing~\citep{siek2006gradual} provided the
theoretical bridge between dynamic and static typing that the Python
type ecosystem
(\texttt{mypy}~\citep{mypy}, \texttt{pyright}~\citep{pyright},
\texttt{pytype}~\citep{pytype}) and TypeScript/Flow have since
operationalized at scale. TensorGuard treats PyTorch tensors as the
gradually-typed fragment: well-annotated functions are checked
precisely, while unknown call sites yield \textsc{Abstain} verdicts
rather than soundness violations or spurious errors.

\subsection{PyTorch- and autograd-specific tooling}

Several recent PyTorch internals tackle adjacent problems.
\texttt{FakeTensorMode}~\citep{faketensor} and
\texttt{torch.export}~\citep{torchexport} symbolically execute graphs
with metadata-only tensors, enabling shape and dtype propagation, but
they (i) require the program to actually run under a tracer and
(ii) do not reason about \texttt{requires\_grad} flow across module
boundaries. \texttt{torch.fx}~\citep{reed2022torchfx} provides the
graph IR these tools build on. TorchDynamo~\citep{ansel2024pytorch2}
inserts runtime guards that check shape/dtype invariants at JIT time;
TensorGuard's contracts can be viewed as the static, refinement-typed
analog of these guards.

The closest comparator is \textbf{Pytea}~\citep{lee2020pytea}, which
performs symbolic shape analysis on PyTorch source via constraint
generation and Z3 discharge. We share Pytea's no-execution stance
and its use of an SMT backend, but differ in three substantive ways:
(1)~Pytea targets shape errors only, whereas TensorGuard's refinement
language additionally tracks \texttt{requires\_grad} and
backward-pass well-formedness; (2)~TensorGuard's checker core is
mechanized in Lean, providing a small trusted base, where Pytea's
constraint translator is unverified Python; (3)~TensorGuard treats
\textsc{Abstain} as a first-class verdict on unhandled idioms,
whereas Pytea reports either ``safe'' or ``unsafe'' and tends to
report many false positives on modern Transformer code.
Hattori et al.~\citep{hattori2023dimension}\,(TODO: verify venue)
introduce dimension-typed APIs for ML libraries; their dimension
algebra is largely orthogonal to ours and could be composed with
TensorGuard's refinements.

\subsection{Lightweight static analysis for ML libraries}

Industrial ML linters such as Pythia~\citep{pythia}\,(TODO: verify
citation), \texttt{pylint}, and \texttt{ruff}~\citep{ruff} catch
syntactic and naming-level bugs but do not reason about tensor
semantics. The runtime-checking layer ---
\texttt{jaxtyping}, \texttt{torchtyping},
\texttt{beartype}~\citep{beartype} --- raises shape errors on the
first offending call but only after the program has been launched
and (typically) GPUs have been allocated. TensorGuard's
contribution at this layer is to push the same class of checks
\emph{ahead} of execution while preserving the lightweight
annotation style these tools popularized.

\begin{table}[t]
 \centering
 \scriptsize
 \setlength{\tabcolsep}{3pt}
 \renewcommand{\arraystretch}{1.15}
 \newcommand{\cmark}{\ensuremath{\checkmark}}
 \newcommand{\xmark}{\ensuremath{\times}}
 \caption{Calibrated comparison along nine PL/ML axes.
 \cmark{} = supported; \xmark{} = not supported;
 $\partial$ = partial / under-developed; --- = not applicable.
 ``Refinement-typed'' means types carry SMT-decidable predicates
 (not just shape literals). ``Symbolic shapes'' means dimensions
 may be unknown variables with constraints. ``Grad-flag tracking''
 means the type system reasons about \texttt{requires\_grad}.
 ``Backward verifier'' means the system checks well-formedness of
 the autograd backward pass, not only the forward computation.
 ``Lean-audited rules'' means the operator shape-transfer
 table is mechanically defined and empirically audited
 against PyTorch (the analyser implementation itself is not
 mechanised).
 ``No-execution'' means the analysis runs without invoking a
 tracer, GPU, or Python interpreter on user code.
 ``Abstention 1st-class'' means \textsc{Abstain} is a reported
 verdict distinct from \textsc{Verified}/\textsc{Refuted}.
 ``Modern Transformer support'' is a coarse practical judgement of
 whether the tool handles HuggingFace-style attention/KV-cache code
 out of the box.}
 \label{tab:contrib}
 \begin{tabular}{l c c c c c c c c c}
 \toprule
 Tool
 & \rotatebox{60}{Refinement-typed}
 & \rotatebox{60}{Symbolic shapes}
 & \rotatebox{60}{Grad-flag tracking}
 & \rotatebox{60}{Backward verifier}
 & \rotatebox{60}{Contracts / A-G}
 & \rotatebox{60}{Lean-audited rules}
 & \rotatebox{60}{No-execution}
 & \rotatebox{60}{Abstention 1st-class}
 & \rotatebox{60}{Modern Transformer} \\
 \midrule
 \textbf{TensorGuard}
 & \cmark & \cmark & \cmark & \cmark & \cmark & \cmark & \cmark & \cmark & \cmark \\
 Pytea
 & $\partial$ & \cmark & \xmark & \xmark & \xmark & \xmark & \cmark & \xmark & $\partial$ \\
 jaxtyping
 & \xmark & $\partial$ & \xmark & \xmark & $\partial$ & \xmark & \xmark & \xmark & \cmark \\
 torchtyping
 & \xmark & $\partial$ & \xmark & \xmark & $\partial$ & \xmark & \xmark & \xmark & \cmark \\
 FakeTensorMode
 & \xmark & \cmark & \xmark & $\partial$ & \xmark & \xmark & \xmark & \xmark & \cmark \\
 torch.export
 & \xmark & \cmark & \xmark & $\partial$ & \xmark & \xmark & \xmark & \xmark & \cmark \\
 MyPy / Pyright
 & \xmark & \xmark & \xmark & \xmark & \xmark & \xmark & \cmark & \xmark & \xmark \\
 beartype
 & \xmark & \xmark & \xmark & \xmark & $\partial$ & \xmark & \xmark & \xmark & \cmark \\
 \bottomrule
 \end{tabular}
\end{table}

